\section{特征提取：从手工设计到自动学习}

\subsection{手工特征提取的局限}

为了正确分类图像，我们需要提取能区分不同类别的特征。传统方法依赖\textbf{手工特征提取}（Manual Feature Extraction），流程如下：

\begin{figure}[H]
\centering
\includegraphics[width=0.9\textwidth]{slides-images/slide-15.jpg}
\caption{手工特征提取流程：领域知识 $\rightarrow$ 定义特征 $\rightarrow$ 检测特征进行分类。核心问题是这一流程对变化极度敏感。\protect\footnotemark}
\end{figure}
\footnotetext{Slides 第15页。}

例如，要识别人脸，你可能定义特征为"鼻子、眼睛、嘴巴"。但这立即引出一个递归问题：如何检测"眼睛"？你又需要定义更底层的特征来描述眼睛的外观。

\begin{warningbox}{手工特征的核心困境}
手工特征提取面临两个根本性挑战：
\begin{enumerate}
    \item \textbf{特征定义的递归性}：定义高层特征需要定义低层特征，而低层特征的定义又需要更低层特征，形成无穷递归。
    \item \textbf{变异性}：同一类物体在不同视角、光照、尺度、形变和遮挡条件下的外观差异巨大，手工规则无法穷举所有情况。
\end{enumerate}
\end{warningbox}

\begin{figure}[H]
\centering
\includegraphics[width=0.9\textwidth]{slides-images/slide-16.jpg}
\caption{图像识别面临的挑战：视角变化、尺度变化、形变、遮挡、光照条件、背景杂乱和类内变化。\protect\footnotemark}
\end{figure}
\footnotetext{Slides 第16页。}

\subsection{学习特征层次结构}

深度学习的核心思想是：\textbf{能否直接从数据中学习特征的层次结构}（hierarchy of features），而不是手工设计？

\begin{figure}[H]
\centering
\includegraphics[width=0.85\textwidth]{slides-images/slide-18.jpg}
\caption{从数据中学习特征层次：低层特征（边缘、暗点）$\rightarrow$ 中层特征（眼睛、耳朵、鼻子）$\rightarrow$ 高层特征（面部结构）。\protect\footnotemark}
\end{figure}
\footnotetext{Slides 第18页。参考 Lee+ ICML 2009。}

\begin{importantbox}{特征层次结构的关键洞察}
视觉特征是分层组织的：
\begin{itemize}
    \item \textbf{低层特征}（Low-level features）：边缘、暗斑等最基本的视觉元素
    \item \textbf{中层特征}（Mid-level features）：由低层特征组合而成，如眼睛、耳朵
    \item \textbf{高层特征}（High-level features）：由中层特征组合而成，如完整的面部结构
\end{itemize}
深度学习的目标就是让网络自动学习这种层次化的特征表示。
\end{importantbox}

\subsection{本章小结}

手工特征提取由于递归性和变异性问题而不可扩展。深度学习的核心贡献是让网络自动从数据中学习分层特征表示，从低层的边缘检测到高层的语义理解，全程无需人工干预。

\newpage
